\subsection{Table and figure handling}\label{subsec:doc-table-and-figure-handling}

Tables and figures are handled together, as the pipeline applies the same operations to both artifact types: it detects the region on the page, masks it out of the body text, crops it into a standalone image, and finally links the crop to a caption. Tables are the harder case and determine the detector design; figures use the same masking, cropping, and caption-linking mechanism.

Some tables on the pages cannot be detected by Docling alone. To address this gap, a Docling and Table Transformer (TATR) hybrid table-detecting architecture is used. The hybrid detector merges the two outputs: it pools the table regions from both detectors, keeping a region if either reports a table, and then merges the resulting boxes geometrically. The merge is iterative and threshold-free: any two boxes with a non-zero overlap are replaced by their bounding union, and the process continues until no overlapping boxes remain. A table reported by both detectors, or split across several overlapping regions, is therefore collapsed into a single region.

The hybrid architecture improves recall, at the cost of potentially lowering precision. Table detections that are almost entirely within the bounding boxes of a figure are dropped, and mislabeled tables at the page headers are dropped to mitigate the precision loss from these fault modes. 

Figure bounding boxes are taken from Docling alone, as TATR is table-specific, and does not affect figure detection. The main failure mode for figures is decorative graphics -- publisher logos and inline icons, mislabeled by Docling as figures. A size filter removes the smallest of these by skipping figure regions below 50 points in height or width. However, moderately-sized decorative graphics are not handled by any other filter, and survive as false positives. These are still masked correctly, and do not contaminate the text output, so the cost is lower precision for figure extraction, not lower quality in the text output. This limitation is documented in Subsection~\ref{subsec:doc-extraction-known-limitations}. 

After tables and figures in a page are located, their bounding boxes are redacted, which removes the underlying text objects rather than merely cover them. The page is then parsed for a second time. This masking step protects the body-text extraction stage: it prevents texts that would be detected in figures and tables from being mistakenly reread as body text and contaminating the text merger input. The masked tables and figures are preserved separately, so nothing the reader needs is lost.

Each detected region is finally cropped to a standalone image. Table crops are extended downward to include footnotes detected nearby, so that table notes are not absorbed into the extracted text. The resulting crops are then matched to a caption using geometry-based heuristics -- the vertical gap between the artifact and caption bounding box edges. The figure or table number is then parsed from the matched caption text. 

Caption linking is purely geometric, and is the least reliable part of the table-and-figure handling stage. It mismatches when tables are rotated, captions are attached to the side of the artifact, or on tables with captions on a different page. The caption-linking and other precision-recall failures are treated as known limitations in Subsection~\ref{subsec:doc-extraction-known-limitations}; the effects on the precision and recall measurements are reported in Chapter~\ref{chapter:Results}.

